\section{Proofs for \texorpdfstring{Section~\ref{sec:abs}}{abstract collectives}}\label{app:abs}

The proofs follow the order of \cref{sec:abs}. Two facts about M-matrices are used throughout. Let $M=D_0-N_0$ be symmetric, with $D_0$ diagonal, $N_0\ge0$ entrywise with zero diagonal, and $M\succeq\delta D_0$ for some $\delta>0$; in finite dimensions the last condition is positive definiteness, and on $\ell^2(\Z)$ we also take $D_0$ bounded above and below by positive constants. Then $S=D_0^{-1/2}N_0D_0^{-1/2}$ is symmetric and nonnegative with $S\preceq(1-\delta)I$, and since $\lvert\langle f,Sf\rangle\rvert\le\langle\lvert f\rvert,S\lvert f\rvert\rangle$, its norm is at most $1-\delta$. Hence $M^{-1}=D_0^{-1/2}\sum_{n\ge0}S^nD_0^{-1/2}=\sum_{n\ge0}(D_0^{-1}N_0)^nD_0^{-1}$, the series converging in norm, so $M^{-1}\ge0$ entrywise, and $M^{-1}$ is positive at every entry when $N_0$ is irreducible, since some term is. If two such operators satisfy $M_1\le M_2$ entrywise, then $M_1^{-1}-M_2^{-1}=M_1^{-1}(M_2-M_1)M_2^{-1}\ge0$.

\begin{proof}[Proof of \cref{prop:abs-decoupling}]
Consider first a level-0 receiver $i$ with any assigned noises. It predicts with $F$ and $Q$ and updates by one observation at a time; observations with independent noises can be processed sequentially, because under the receiver's own model the sequential and the joint Gaussian updates coincide. Suppose $P_i$ is block diagonal across contents. The prediction $FP_iF^\top+Q$ is block diagonal by (a). An observation whose model matrix $\hat H$ has nonzero columns only in block $c$, a sensing channel by (c) or a modality-$c$ message by block purity, has gain $\Kg=P^-_i\hat H^\top(\hat HP^-_i\hat H^\top+\hat R)^{-1}$ with nonzero rows only in block $c$, and $(I-\Kg\hat H)P^-_i$ differs from $P^-_i$ only in the $(c,c)$ block. By induction over rounds every posterior covariance and every gain is block diagonal, and under (d) the gains are deterministic. The error obeys $e^-_i=Fe_i-w_t$ in the prediction and $e^+_i=e^-_i+\Kg(o-\hat Hm^-_i)$ in an update. For a modality-$c$ message $o=H_jm_j+\nu=H_j(\eta+e_j)+\nu$, where $H_j$ reads only the $c$-block of $m_j$, so $o-\hat Hm^-_i=(H_j-\hat H)\eta+H_je_j-\hat He^-_i+\nu$ involves only the $c$-blocks of $\eta$, $e_j$ and $e^-_i$, and $\Kg$ writes only into block $c$; sensing updates are alike. Hence $e^{(c)}(t+1)$ is a linear function, with deterministic coefficients, of $e^{(c)}(t)$, of $\eta^{(c)}_t$ and $w^{(c)}_t$, and of the noises of content-$c$ channels, and the coefficient of $e^{(c)}_j$ in the update of $e^{(c)}_i$ vanishes unless $i=j$ or $i$ reads $j$ in modality $c$, an edge of $\mathcal G_c$. This is the recursion with transfer $W^{(c)}(t)$. The world blocks evolve independently by (a), the noises of different contents are independent by (c), and so are the initial quantities by (b); the processes $(e^{(c)},\eta^{(c)})$, $c\in\mathcal M$, are therefore independent.

For exact fusion we prove by induction over rounds that every agent's estimate of $\eta^{(c)}_t$ is a linear functional of the content-$c$ primitives: the path of $\eta^{(c)}$, the noises of content-$c$ channels and the content-$c$ initial quantities. If this holds up to round $t$, every content-$c$ observation at round $t$, sensing or message, is such a functional. The primitives of different contents are independent, so the conditional law of $\eta^{(c)}_t$ given all of an agent's data equals its conditional law given the agent's content-$c$ data, which proves the claim at the next round. Posterior covariances are then block diagonal, and the error processes of different contents are independent. The results of \cref{sec:sca} cited in the statement assume a linear error recursion with local noise on a fixed schedule, which each $e^{(c)}$ satisfies at level 0.
\end{proof}

\begin{proof}[Proof of \cref{prop:abs-crossover}]
The update is symmetric and translation invariant with symbol $\hat W(k)=\bigl(\begin{smallmatrix}c_1+2a_1\cos k&a_\times\\a_\times&c_2+2a_2\cos k\end{smallmatrix}\bigr)$. Since $\hat W(0)-\hat W(k)$ and $\hat W(k)-\hat W(\pi)$ are diagonal and positive semidefinite, every eigenvalue of $\hat W(k)$ lies between the smallest eigenvalue of $\hat W(\pi)$ and the largest of $\hat W(0)$. Now $I-\hat W(0)=\bigl(\begin{smallmatrix}\varepsilon_1&-a_\times\\-a_\times&\varepsilon_2\end{smallmatrix}\bigr)\succ0$ because $a_\times^2<\varepsilon_1\varepsilon_2$, and $I+\hat W(\pi)=\bigl(\begin{smallmatrix}2c_1+\varepsilon_1&a_\times\\a_\times&2c_2+\varepsilon_2\end{smallmatrix}\bigr)\succ0$ because $(2c_1+\varepsilon_1)(2c_2+\varepsilon_2)\ge\varepsilon_1\varepsilon_2>a_\times^2$. Hence $\rho(W)<1$, and since $W$ is symmetric and the noise is $qI$, the stationary covariance is $\Sigma=q\sum_{n\ge0}W^{2n}=q(I-W^2)^{-1}=\frac q2[(I-W)^{-1}+(I+W)^{-1}]$.

For $u=\cosh m>1$, two geometric series give $\sum_{d\in\Z}e^{-m\lvert d\rvert}e^{ikd}=\sinh m/(u-\cos k)$, so the Fourier coefficients of $1/(u-\cos k)$ are $e^{-m\lvert d\rvert}/\sinh m$, and, replacing $k$ by $k+\pi$, those of $1/(u+\cos k)$ are $(-1)^de^{-m\lvert d\rvert}/\sinh m$. With $z=\cos k$, $[(I-\hat W)^{-1}]_{11}=A_2(z)/(A_1(z)A_2(z)-a_\times^2)$. Since $A_c(z)=2a_c(\cosh m^0_c-z)$, the denominator is $4a_1a_2(z-u_-)(z-u_+)$, and because $A_2$ is affine, partial fractions give $\sum_\pm A_2(u_\pm)/[4a_1a_2(u_\pm-u_\mp)(z-u_\pm)]$, where $u_+\ne u_-$ because $\cosh m^0_1\ne\cosh m^0_2$. As $-1/(z-u_\pm)=1/(u_\pm-\cos k)$, this gives the non-alternating terms with the stated $\mathcal A_\pm$. In the same way $[(I+\hat W)^{-1}]_{11}=B_2(z)/(B_1(z)B_2(z)-a_\times^2)=\sum_\pm B_2(w_\pm)/[4a_1a_2(w_\pm-w_\mp)(z-w_\pm)]$, with $w_+\ne w_-$ when $a_\times\ne0$, and $1/(z-w_\pm)=1/(\cosh m'_\pm+\cos k)$ gives the alternating terms.

(a) The roots of $A_1A_2=a_\times^2$ solve $(u-\cosh m^0_1)(u-\cosh m^0_2)=a_\times^2/(4a_1a_2)\ge0$, and the left side is negative strictly between $\cosh m^0_2$ and $\cosh m^0_1$, so $u_-\le\cosh m^0_2$ and $u_+\ge\cosh m^0_1$, strictly when $a_\times\ne0$. At $u=1$, $A_1A_2-a_\times^2=\varepsilon_1\varepsilon_2-a_\times^2>0$, so $1$ lies outside $[u_-,u_+]$, an interval containing $[\cosh m^0_2,\cosh m^0_1]\subset(1,\infty)$; hence $u_->1$. Similarly $B_c(w)=2a_c(w-w^0_c)$ with $w^0_c=-(1+c_c)/(2a_c)<-1$, because $1+c_c-2a_c=2c_c+\varepsilon_c>0$; the roots $w_\pm$ bracket the $w^0_c$, and $B_1(-1)B_2(-1)-a_\times^2=(2c_1+\varepsilon_1)(2c_2+\varepsilon_2)-a_\times^2>0$ places $-1$ outside $[w_-,w_+]$, so $w_+<-1$.

(b) $A_2(u_+)=2a_2(\cosh m^0_2-u_+)<0$ and $u_+-u_->0$ give $\mathcal A_+>0$, and for $a_\times\ne0$, $A_2(u_-)=2a_2(\cosh m^0_2-u_-)>0$ and $u_--u_+<0$ give $\mathcal A_->0$. For $u\in[1,u_-]$, $A_c(u)\ge A_c(u_-)>0$, and $A_1(u)A_2(u)\ge a_\times^2$ because the quadratic $A_1A_2-a_\times^2$ is nonnegative to the left of its smaller root. Since $B_c(-u)=A_c(u)+2c_c$ and $c_1,c_2\ge0$, if $c_1+c_2>0$ then every $u\in[1,u_-]$ satisfies
\[
B_1(-u)B_2(-u)-a_\times^2\ge2c_1A_2(u)+2c_2A_1(u)+4c_1c_2>0.
\]
So the quadratic $B_1B_2-a_\times^2$, which is negative strictly between its roots, is positive on $[-u_-,-1]$; both roots lie below $-1$, hence below $-u_-$, and $\cosh m'_\pm=-w_\pm>u_-=\cosh m_-$. The alternating terms and the term in $m_+>m_-$ therefore decay faster than $e^{-m_-\lvert d\rvert}$.

(c) Evaluating the quadratic at $u_-$ gives
\[
\cosh m^0_2-u_-=\frac{a_\times^2}{4a_1a_2(\cosh m^0_1-u_-)}=\frac{a_\times^2}{4a_1a_2(\cosh m^0_1-\cosh m^0_2)}+O(a_\times^4),
\]
and $u_+-\cosh m^0_1=O(a_\times^2)$. Hence $A_2(u_-)=a_\times^2/[2a_1(\cosh m^0_1-\cosh m^0_2)]+O(a_\times^4)$, $A_2(u_+)=-2a_2(\cosh m^0_1-\cosh m^0_2)+O(a_\times^2)$, $u_+-u_-=\cosh m^0_1-\cosh m^0_2+O(a_\times^2)$ and $\sinh m_\pm=\sinh m^0_{1,2}+O(a_\times^2)$. So $\mathcal A_+=1/(2a_1\sinh m^0_1)+O(a_\times^2)$ and $\mathcal A_-=a_\times^2/[8a_1^2a_2(\cosh m^0_1-\cosh m^0_2)^2\sinh m^0_2]+O(a_\times^4)$, whose ratio is the stated one. The two exponentials are equal where $\mathcal A_+e^{-m_+d}=\mathcal A_-e^{-m_-d}$, which is the first form of $d_\times$. Inserting the ratio and $m_+-m_-=m^0_1-m^0_2+O(a_\times^2)$ gives the second, the error being the product of the logarithm with the correction to the denominator.
\end{proof}

\begin{proof}[Proof of \cref{prop:abs-homogeneous}]
For distinct $i,j$ every chain from $i$ to $j$ contains a link of weight at least $\varpi^\circ$, so $u(i,j)\ge\varpi^\circ$, and $u_{\mathrm{in}}(C)\ge\varpi^\circ$ whenever $\lvert C\rvert\ge2$. Since $\mathcal G_c$ is connected and $C$ is a proper nonempty subset, some edge of $\mathcal G_c$ joins $C$ to its complement, and its weight $\varpi^\circ$ bounds $\varpi_{\mathrm{out}}(C)$ from above. Hence $\mathcal U(C)\le0$.
\end{proof}

\begin{proof}[Proof of \cref{prop:abs-barrier}]
(a) Read $\varepsilon I+L_b$ as the admittance matrix of a resistor network with conductance $a$ on every bond except $a_b$ on $(0,1)$, and conductance $\varepsilon$ from every node to ground; $\mathsf G(x,y)$ is the potential at $x$ when unit current enters at $y$, decaying at infinity. A semi-infinite uniform chain with first node 1 carries the decaying potentials $\phi_n=\phi_1e^{-m(n-1)}$, since $\varepsilon\phi_n+a(2\phi_n-\phi_{n-1}-\phi_{n+1})=0$ for $n\ge2$ requires $e^m+e^{-m}=2+\varepsilon/a$. The current entering node 1 is then $\phi_1[\varepsilon+a(1-e^{-m})]=\phi_1a(e^m-1)$, so the input admittance is $Y=a(e^m-1)$. When unit current enters at node 1, node 1 heads the right half-chain and is joined through $a_b$ to node 0, which heads the left half-chain; the admittance seen at node 1 is $Y+a_bY/(a_b+Y)$, so $\mathsf G(1,1)=(a_b+Y)/[Y(2a_b+Y)]$, and the voltage divider gives $\mathsf G(0,1)=\mathsf G(1,1)a_b/(a_b+Y)=a_b/[Y(2a_b+Y)]$. No current enters at $x\le0$, so $\mathsf G(x,1)=\mathsf G(0,1)e^{-m\lvert x\rvert}$, and by the symmetry of $\mathsf G$ and the same argument on the right, $\mathsf G(x,y)=\mathsf G(0,1)e^{-m(\lvert x\rvert+y-1)}$ for $x\le0<y$. For $n_b=1$ this gives $\mathsf G^0(x,y)=e^{-m(y-x)}/(2a\sinh m)$. The ratio is $2ae^m\mathsf G(0,1)\sinh m$, and with $e^m(e^m-e^{-m})=(e^m-1)(e^m+1)$ it equals $a_b(e^m+1)/(2a_b+a(e^m-1))=(e^m+1)/(2+n_b(e^m-1))$. Dividing numerator and denominator by $e^m-1$ gives the second form.

(b) With $\mathfrak s=e^m-1=m+O(m^2)$, $\mathcal T=(1+\mathfrak s/2)/(1+n_b\mathfrak s/2)=1-(n_b-1)\mathfrak s/2+O(n_b^2\mathfrak s^2)=1-\frac m2(n_b-1)+O(m^2n_b^2)$. From the second form, $\mathcal Tn_bm/2=\frac m2\coth(m/2)\cdot n_b/(n_b+2/(e^m-1))$; the first factor tends to one as $m\to0$, and the second tends to one when $mn_b\to\infty$, because $2/(e^m-1)\le2/m$. Taking logarithms gives the form of $d_b$.

(c) The update $W=(1-\varepsilon)I-L_b$ is symmetric with spectrum in $[1-\varepsilon-4a,1-\varepsilon]$, and $1-\varepsilon-4a=2c+\varepsilon-1>-1$, so $\rho(W)<1$ and $\Sigma=\frac q2[\mathsf G+(I+W)^{-1}]$, and likewise for $\Sigma^0$. Let $D_\pm=\diag((-1)^x)$ and let $\mathrm{Deg}_b$ be the diagonal matrix of weighted degrees. The chain is bipartite, so $D_\pm L_bD_\pm=2\mathrm{Deg}_b-L_b$ and $\tilde M:=D_\pm(I+W)D_\pm=(2-\varepsilon)I-2\mathrm{Deg}_b+L_b$, whose diagonal is $2-\varepsilon-2a=2c+\varepsilon+2a$ in the bulk and $2-\varepsilon-a-a_b$, larger by $a-a_b\ge0$, at the sites 0 and 1, and whose off-diagonal entries are $-a$, or $-a_b\ge-a$ on the barrier bond. The uniform operator $\tilde M^0=(2c+\varepsilon)I+L_0$ therefore satisfies $\tilde M^0\le\tilde M$ entrywise. Both are of the form of the preliminary facts, with $\tilde M\succeq(2c+\varepsilon)I$ because $\tilde M$ is unitarily equivalent to $I+W$, and likewise $\tilde M^0$, so $0\le\tilde M^{-1}\le(\tilde M^0)^{-1}$ entrywise. The symbol of $\tilde M^0$ is $2c+\varepsilon+2a(1-\cos k)=2a(\cosh m'-\cos k)$ with $\cosh m'=1+(2c+\varepsilon)/(2a)=(1+c)/(2a)$, using $2a+c+\varepsilon=1$, so $(\tilde M^0)^{-1}(x,y)=e^{-m'\lvert x-y\rvert}/(2a\sinh m')$ by the Fourier fact in the previous proof. Since $(I+W)^{-1}(x,y)=(-1)^{x-y}\tilde M^{-1}(x,y)$, we get $\lvert(I+W)^{-1}(x,y)\rvert\le e^{-m'(y-x)}/(2a\sinh m')$, with equality, and the sign $(-1)^{x-y}$, for the uniform chain. Also $m'>m$, because $(1+c)/(2a)>(1-c)/(2a)=\cosh m$ when $c>0$. Write $g=\mathsf G^0(x,y)>0$, so that $e^{-m'(y-x)}/(2a\sinh m')=r_{xy}g$. By (a), $\Sigma(x,y)=\frac q2g(\mathcal T+\mathfrak e_1)$ and $\Sigma^0(x,y)=\frac q2g(1+\mathfrak e_0)$ with $\lvert\mathfrak e_0\rvert,\lvert\mathfrak e_1\rvert\le r_{xy}$, hence
\[
\Bigl\lvert\frac{\Sigma(x,y)}{\Sigma^0(x,y)}-\mathcal T\Bigr\rvert=\frac{\lvert\mathfrak e_1-\mathcal T\mathfrak e_0\rvert}{\lvert1+\mathfrak e_0\rvert}\le\frac{(1+\mathcal T)r_{xy}}{1-r_{xy}}\le\frac{2r_{xy}}{1-r_{xy}},
\]
since $\mathcal T\le1$ for $n_b\ge1$.
\end{proof}

\begin{proof}[Proof of \cref{prop:abs-bypass}]
The matrix $I-W$ is positive definite, because $A_1\succeq\varepsilon_1I$, $A_2\succeq\varepsilon_2I$ and its Schur complement satisfies $A_1-a_\times^2A_2^{-1}\succeq(\varepsilon_1-a_\times^2/\varepsilon_2)I\succ0$. Block inversion gives $\mathsf G_{11}=(A_1-a_\times^2A_2^{-1})^{-1}=\sum_{n\ge0}(a_\times^2A_1^{-1}A_2^{-1})^nA_1^{-1}$, where the series converges because $\lVert a_\times^2A_1^{-1}A_2^{-1}\rVert_2\le a_\times^2/(\varepsilon_1\varepsilon_2)<1$. Each $A_c$ is positive definite with nonpositive off-diagonal entries, so $A_c^{-1}\ge0$ entrywise and every term is nonnegative. Since $L_1$ has no edge across $C$, $A_1$ and $A_1^{-1}$ are block diagonal with respect to $(C,V\setminus C)$, so the term $n=0$ vanishes across the cut, which gives the identity and the bounds. If $A_2$ is also block diagonal, every term is, and $\mathsf G_{11}$ vanishes across the cut. If $L_2$ is connected, $A_2$ is an irreducible nonsingular M-matrix and $A_2^{-1}$ is entrywise positive, and the diagonal entries of the positive definite matrix $A_1^{-1}$ are positive; since all terms of the sum are nonnegative, $[A_1^{-1}A_2^{-1}A_1^{-1}](x,y)=\sum_{x'\in C,y'\notin C}A_1^{-1}(x,x')A_2^{-1}(x',y')A_1^{-1}(y',y)\ge A_1^{-1}(x,x)A_2^{-1}(x,y)A_1^{-1}(y,y)>0$, and the lower bound $a_\times^2[A_1^{-1}A_2^{-1}A_1^{-1}](x,y)$ is positive when $a_\times\ne0$. At $a_\times=0$ it vanishes, as does $\mathsf G_{11}(x,y)$.
\end{proof}

\begin{proof}[Proof of \cref{prop:abs-conformity}]
(a) $\ln[q_{\mathrm p}/(1-q_{\mathrm p})]=b_{\mathrm p}\ln[x/(1-x)]$, so $q_{\mathrm p}(\frac12)=\frac12$ and, since $\operatorname{sig}'(0)=\frac14$, the slope of $q_{\mathrm p}$ at $\frac12$ is $b_{\mathrm p}\cdot\frac14\cdot4=b_{\mathrm p}>1$, which puts $q_{\mathrm p}$ above the diagonal immediately to the right of $\frac12$. Further $q_{\mathrm p}(0)=0$ for $b_{\mathrm p}>0$, while $q(0)=\operatorname{sig}(-\gamma J\mathsf d)>0$ because $\operatorname{sig}$ is positive.

(b) Both rules equal $\frac12$ at $x=\frac12$, and the slope of $q$ there is $2\gamma J\mathsf d\operatorname{sig}'(0)=\gamma J\mathsf d/2$, so the slopes agree exactly when $b_{\mathrm p}=\gamma J\mathsf d/2$. Then the log-odds of $q$, $2b_{\mathrm p}(2x-1)$, is the tangent at $\frac12$ of the log-odds $\ell_{\mathrm p}(x)=b_{\mathrm p}\ln[x/(1-x)]$ of $q_{\mathrm p}$. Since $\ell_{\mathrm p}''(x)=b_{\mathrm p}[(1-x)^{-2}-x^{-2}]$ is negative on $(0,\frac12)$ and positive on $(\frac12,1)$, $\ell_{\mathrm p}$ lies strictly below its tangent on $(0,\frac12)$ and strictly above it on $(\frac12,1)$, and $\operatorname{sig}$ is increasing.

(c) An end agent of the chain has one partner, so its local frequency is 0 or 1 and the power rule copies its partner; an interior agent has $x\in\{0,\frac12,1\}$ and adopts $+1$ with probability $0$, $\frac12$ or $1$. The two uniform configurations are therefore absorbing, and no revision increases the number of walls. In any other configuration take the leftmost wall, between the agents $i$ and $i+1$. If $i$ is the end agent, its revision removes the wall and creates none; otherwise agent $i$ sees partners of both codes and with probability $\frac12$ switches, which moves the wall one bond to the left and creates no other wall. From every configuration with walls, a finite sequence of revisions, each of positive probability, therefore lowers the number of walls by one, and the finite chain is absorbed in a uniform configuration with probability one. Under $q$ at finite $\gamma$ every change of a single code has positive probability, since $\operatorname{sig}$ takes values in $(0,1)$, so the chain is irreducible. With values affine in the number of matched partners and a slope $J$ common to all agents, as here, the code game is an exact potential game on every graph (\cref{prop:coh-code-potential}(c)), here the pairwise game with coupling $J$ on every edge of the chain, and $q$ is its logit revision with a uniformly drawn reviser, so the stationary law is the Gibbs law of \cref{prop:coh-code-potential}(a), described in \cref{prop:abs-seguy}(c).
\end{proof}

\begin{proof}[Proof of \cref{prop:abs-seguy}]
(a) Write $X=\langle Z,\mathsf u_1\rangle$ and $Y=\langle Z,\mathsf u_2\rangle$ for an isotropic standard Gaussian vector $Z\in\R^2$ and unit vectors with $\langle\mathsf u_1,\mathsf u_2\rangle=\Gamma$. The signs of $X$ and $Y$ differ exactly when the line orthogonal to $Z$ separates $\mathsf u_1$ from $\mathsf u_2$; the direction of $Z$ is uniform, so the probability is the angle $\arccos\Gamma$ divided by $\pi$, which is Sheppard's formula \citep{sheppard1899application}. Linearity of expectation gives $\mathcal S(d)$, and the strong law of large numbers over independent features gives almost sure convergence. The expansions follow from $\arccos(1-\mathsf h)=(2\mathsf h)^{1/2}(1+O(\mathsf h))$ as $\mathsf h\to0$, with $1-e^{-d/\xi}=(d/\xi)(1+O(d/\xi))$ and $1-e^{-d^2/(2\sigma_{\mathrm g}^2)}=(d^2/2\sigma_{\mathrm g}^2)(1+O(d^2/\sigma_{\mathrm g}^2))$, and from $\arccos\Gamma=\frac\pi2-\arcsin\Gamma$ with $\arcsin\Gamma=\Gamma(1+O(\Gamma^2))$.

(b) $\sigma_0\ne\sigma_d$ exactly when $n_{\mathrm w}(d)$ is odd, and every integer $n\ge0$ satisfies $n-n(n-1)\le\one[n\text{ odd}]\le n$: both sides equal $n$ for $n\in\{0,1\}$, and for $n\ge2$ the left side is $n(2-n)\le0$. Taking expectations and using $\E n_{\mathrm w}(d)=dp_{\mathrm w}$, which follows from stationarity, gives the bounds; the linear rise follows when the second term is $o(dp_{\mathrm w})$.

(c) By \cref{prop:coh-code-potential}(a), asynchronous logit revision of the pairwise game with coupling $J$ on the edges of the chain is reversible, and at $h_i=0$ its stationary law is proportional to $\exp(-\gamma J\sum_x\one[\sigma_x\ne\sigma_{x+1}])$. On an open chain the map from configurations to the first code and the bond variables $\sigma_x\sigma_{x+1}$ is a bijection, so the bond variables are independent, each equal to $-1$ with probability $e^{-\gamma J}/(1+e^{-\gamma J})=1/(1+e^{\gamma J})$. Then $\E[\sigma_0\sigma_d]=\bigl(1-2/(1+e^{\gamma J})\bigr)^d=\tanh^d(\gamma J/2)$, and $\Pr(\sigma_0\ne\sigma_d)=\frac12(1-\E[\sigma_0\sigma_d])$.

(d) For a standard Gaussian pair with correlation $\Gamma$, $\E[\operatorname{sign}X\operatorname{sign}Y]=1-2\arccos(\Gamma)/\pi=\frac2\pi\arcsin\Gamma$, which is strictly increasing in $\Gamma$ and nonnegative for $\Gamma\ge0$. The readout weights $-\ln(\frac2\pi\arcsin\Gamma_{ij})$ are therefore a strictly increasing function of the weights $\varpi_{ij}=-\ln\Gamma_{ij}$ of the field. The hypotheses of \cref{prop:abs-homogeneous}, a common value on the edges and nothing smaller elsewhere, are preserved by a strictly increasing transformation, and the conclusion follows from that proposition.
\end{proof}

\begin{proof}[Proof of \cref{prop:abs-exchange-consistency}]
(a) A potential makes the sums around cycles telescope to zero. Conversely, if every cycle sum vanishes, fix a root and define $\bar\alpha$ at each node by summing $\alpha$ along any path from the root; the vanishing cycle sums make the value independent of the path, and two potentials differ by a constant on the connected graph. The least-squares residual is zero exactly when $\alpha$ lies in the range of $B$, that is, when a potential exists.

(b) $B\bar\alpha$ has zero holonomy by (a), so only $\phi$ contributes; the least-squares residual differs from $\alpha$ by an element of the range of $B$, which has zero holonomy. The holonomy is a signed sum of independent terms of mean zero, so its second moment is the sum of their variances.

(c) The least-squares estimate through the origin is $\Theta_{ij}=\sum_t\eta_ty_t/\sum_t\eta_t^2=\Theta^\ast_{ij}+\sum_t\eta_t\nu_t/\sum_t\eta_t^2$, where $y_t$ are the messages; given the $\eta_t$ it is Gaussian with mean $\Theta^\ast_{ij}$ and variance $\sigma_y^2/\sum_t\eta_t^2$. Given the $\eta_t$, the law of $n_{ij}^{1/2}(\Theta_{ij}-\Theta^\ast_{ij})$ is therefore exactly $\Ncal(0,\sigma_y^2n_{ij}/\sum_t\eta_t^2)$. The almost sure convergence $n_{ij}^{-1}\sum_t\eta_t^2\to\sigma_\eta^2$, which for independent and identically distributed $\eta_t$ is the strong law of large numbers and for a stationary ergodic sequence is the ergodic theorem, makes this conditional law converge to $\Ncal(0,\sigma_y^2/\sigma_\eta^2)$ for almost every sequence of world values, and by dominated convergence of characteristic functions so does the unconditional law. The event $\Theta_{ij}>0$, on which the logarithm is defined, has probability tending to one, and the delta method with the logarithm, differentiable at $\Theta^\ast_{ij}>0$, gives the stated limit.
\end{proof}

\begin{proof}[Proof of \cref{prop:abs-walras}]
(a) With only its own signal, agent $i$ has $\E_i\eta=(\bar\eta_0/P_0+s_i/R^s_i)/(1/P_0+1/R^s_i)$ and $\Var_i\eta=1/(1/P_0+1/R^s_i)$, so its demand is $A_{\mathrm{ra}}^{-1}[\bar\eta_0/P_0+s_i/R^s_i-(1/P_0+1/R^s_i)p]$. Setting the sum of demands to zero gives $p^\circ$. With $R^s_i=R^s$, $p^\circ=(\bar\eta_0/P_0+\bar s/R^s)/(1/P_0+1/R^s)$ with $\bar s=N^{-1}\sum_is_i$, and $\bar s\to\eta$ almost surely by the strong law of large numbers, conditionally on $\eta$.

(b) By Gaussian conjugacy the pooled posterior is $\Ncal(\hat\eta_{\mathrm{pool}},1/\Lambda_{\mathrm{pool}})$, and it depends on the signals only through $\sum_js_j/R^s_j$. If $p$ is an invertible function of this statistic, the tower property gives $\E[\eta\mid s_i,p]=\E\bigl[\E[\eta\mid s_1,\dots,s_N]\bigm|s_i,p\bigr]=\hat\eta_{\mathrm{pool}}$, and the same argument applies to the whole posterior law. If every agent conditions on $p=\hat\eta_{\mathrm{pool}}$, every agent holds the pooled posterior and demands $(\hat\eta_{\mathrm{pool}}-p)\Lambda_{\mathrm{pool}}/A_{\mathrm{ra}}=0$, so the market clears.

(c) The level-0 posterior is the Gaussian update of the prior $\Ncal(\bar\eta_0,P_0)$ by $s_i$ and $p$ treated as independent observations with precisions $1/R^s_i$ and $h_p$, which has the stated mean and precision. The price errs by $\hat\eta_{\mathrm{pool}}-\eta$, of variance $1/\Lambda_{\mathrm{pool}}$, and by (b) the true posterior given $(s_i,p)$ has precision $\Lambda_{\mathrm{pool}}$, so the excess of the believed precision at $h_p=\Lambda_{\mathrm{pool}}$ is $1/P_0+1/R^s_i$.
\end{proof}

\begin{proof}[Proof of \cref{prop:abs-hayek}]
If $\psi_i=M^\top\mathsf c_1$ for some vector $\mathsf c_1$, then $\psi_i^\top\eta=\mathsf c_1^\top M\eta$ is a function of the conditioning variables and its conditional variance is zero. Conversely, for jointly Gaussian variables a conditional variance vanishes exactly when the variable is almost surely an affine function of the conditioning variables, so $\psi_i^\top\eta=\mathsf c_1^\top M\eta+\mathsf c_2^\top(H^s_i\eta+v_i)+\mathsf c_0$ almost surely for some vectors $\mathsf c_1$, $\mathsf c_2$ and a constant $\mathsf c_0$. Taking the covariance of both sides with $v_i$, which is independent of $\eta$, gives $0=R^s_i\mathsf c_2$, so $\mathsf c_2=0$. Then $(\psi_i-M^\top\mathsf c_1)^\top\eta$ is almost surely constant, and $P_0\succ0$ forces $\psi_i=M^\top\mathsf c_1$. A public message suffices for every agent if and only if every $\psi_i$ lies in $\operatorname{row}(M)$, so $\rank M\ge\dim\operatorname{span}\{\psi_i\}$, with equality for the matrix whose rows are a basis of $\operatorname{span}\{\psi_i\}$.
\end{proof}

\begin{proof}[Proof of \cref{prop:abs-polity}]
(a) Under rule C, stationarity gives $\E(\varphi_x-\varphi_{x_0})^2=2\sigma_\varphi^2(1-\Gamma(x-x_0))$, so a member at $x$ pays $g(1-\Gamma(x-x_0))$; summing over $C$, adding $f$ and dividing by $\lvert C\rvert$ gives $\mathcal C_{\mathrm C}$. Since $\Gamma\ge0$, every finite sum $\sum_{x\in C}\Gamma(x-x_0)$ is at most $\chi$, so $\mathcal C_{\mathrm C}\to g$ whenever $\lvert C\rvert\to\infty$. Further, $\mathcal C_{\mathrm C}(C)<g$ if and only if $\sum_{x\in C}\Gamma(x-x_0)>f/g$; the supremum of these sums over finite $C\ni x_0$ is $\chi$, so such a $C$ exists if $f<g\chi$, and none exists if $f\ge g\chi$. At fixed size, $\mathcal C_{\mathrm C}$ decreases in $\sum_{x\in C}\Gamma(x-x_0)$, which the sites of largest $\Gamma(\cdot-x_0)$ maximize.

(b) Under rule M the policy is $\bar\varphi_C=\lvert C\rvert^{-1}\sum_{x\in C}\varphi_x$, and with $\mu_\varphi=\E\varphi_x$,
\[
\sum_{x\in C}\E(\varphi_x-\bar\varphi_C)^2=\sum_{x\in C}\E(\varphi_x-\mu_\varphi)^2-\lvert C\rvert\Var(\bar\varphi_C)=\lvert C\rvert\sigma_\varphi^2-\frac{\sigma_\varphi^2}{\lvert C\rvert}\sum_{x,y\in C}\Gamma(x-y).
\]
Multiplying by $g/(2\sigma_\varphi^2)$, adding $f$ and dividing by $\lvert C\rvert$ gives $\mathcal C_{\mathrm M}$. For each $x$, $\sum_{y\in C}\Gamma(x-y)\le\chi$, so $\chi_C\le\chi$ and $\mathcal C_{\mathrm M}\to g/2$ whenever $\lvert C\rvert\to\infty$. Further, $\mathcal C_{\mathrm M}(C)<g/2$ if and only if $\chi_C>2f/g$. For boxes $C_n$ of side $n$, $\chi_{C_n}=\sum_z\Gamma(z)\lvert C_n\cap(C_n-z)\rvert/\lvert C_n\rvert\to\chi$ by dominated convergence, since each ratio of cardinalities tends to one; so if $f<g\chi/2$ some box has $\chi_C>2f/g$, and if $f\ge g\chi/2$ no finite $C$ does.

(c) Rule C: for $C=[-d/2,d/2]$ with $x_0=0$, $\int_C\Gamma=2\xi(1-e^{-d/2\xi})$, so with $\varrho=d/(2\xi)$, $(\mathcal C_{\mathrm C}-g)/g=[\bar f_{\mathrm C}/2-1+e^{-\varrho}]/\varrho$. Its derivative in $\varrho$ has the sign of $1-\bar f_{\mathrm C}/2-e^{-\varrho}(1+\varrho)$, and $e^{-\varrho}(1+\varrho)$ decreases from 1 to 0 on $(0,\infty)$ (its derivative is $-\varrho e^{-\varrho}$), so the derivative vanishes at exactly one point $\varrho_{\mathrm C}$ when $0<\bar f_{\mathrm C}<2$. There the bracket equals $-\varrho_{\mathrm C}e^{-\varrho_{\mathrm C}}<0$; the bracket divided by $\varrho$ tends to $+\infty$ as $\varrho\to0$ and to $0$ from below as $\varrho\to\infty$, so $\varrho_{\mathrm C}$ is the global minimizer and the minimum lies below $g$. If $\bar f_{\mathrm C}\ge2$ the bracket is positive and the derivative negative for every $\varrho$. The expansion $e^{-\varrho}(1+\varrho)=1-\varrho^2/2+O(\varrho^3)$ gives $\varrho_{\mathrm C}=\bar f_{\mathrm C}^{1/2}(1+O(\bar f_{\mathrm C}^{1/2}))$, that is $d^\ast_{\mathrm C}=2(f\xi/g)^{1/2}(1+O(\bar f_{\mathrm C}^{1/2}))$, and the root equation can be written $\varrho_{\mathrm C}-\ln(1+\varrho_{\mathrm C})=\ln[1/(1-\bar f_{\mathrm C}/2)]$, so $\varrho_{\mathrm C}$ diverges like the logarithm of $1/(1-\bar f_{\mathrm C}/2)$.

Rule M: $\int_0^d\int_0^de^{-\lvert x-y\rvert/\xi}\mathrm dx\mathrm dy=2\xi d-2\xi^2(1-e^{-d/\xi})$, so with $\varrho=d/\xi$, $2\mathcal C_{\mathrm M}/g=1+[\bar f_{\mathrm M}-2+2(1-e^{-\varrho})/\varrho]/\varrho$. The derivative of the bracket divided by $\varrho$ vanishes exactly when $2-\bar f_{\mathrm M}=\psi_{\mathrm M}(\varrho)$. Now $\varrho^2\psi_{\mathrm M}'(\varrho)=e^{-\varrho}(4+4\varrho+2\varrho^2)-4$, which vanishes at $\varrho=0$ and has derivative $-2\varrho^2e^{-\varrho}<0$, so $\psi_{\mathrm M}$ decreases strictly from 2 to 0 on $(0,\infty)$ and the root $\varrho_{\mathrm M}$ is unique for $0<\bar f_{\mathrm M}<2$. The bracket divided by $\varrho$ tends to $+\infty$ as $\varrho\to0$ and to $0$ from below as $\varrho\to\infty$, so the stationary point is the global minimizer and the minimum lies below $g/2$; if $\bar f_{\mathrm M}\ge2$ the bracket is positive for all $\varrho$. The expansions $\psi_{\mathrm M}(\varrho)=2-\varrho^2/3+O(\varrho^3)$ and $\psi_{\mathrm M}(\varrho)=4/\varrho+O(e^{-\varrho})$ give $\varrho_{\mathrm M}=(3\bar f_{\mathrm M})^{1/2}(1+O(\bar f_{\mathrm M}^{1/2}))$ and $\varrho_{\mathrm M}=4(2-\bar f_{\mathrm M})^{-1}(1+o(1))$, which are the stated limits.

(d) Put $\bar f=f/(g\xi)\in(0,1)$, so that $\bar f_{\mathrm C}=\bar f$ and $\bar f_{\mathrm M}=2\bar f<2$, and write $\varrho=\varrho_{\mathrm C}$. The root equation of rule C gives $2-\bar f_{\mathrm M}=2-2\bar f=4e^{-\varrho}(1+\varrho)-2$. Since $\psi_{\mathrm M}$ is strictly decreasing, $d^\ast_{\mathrm M}>d^\ast_{\mathrm C}$, that is $\varrho_{\mathrm M}>2\varrho$, holds if and only if $\psi_{\mathrm M}(2\varrho)>4e^{-\varrho}(1+\varrho)-2$. Multiplying the difference of the two sides by $\varrho e^{2\varrho}/2>0$ gives
\[
\varrho e^{2\varrho}+e^{2\varrho}-1-\varrho-2\varrho(1+\varrho)e^\varrho=\sum_{\mathsf n\ge3}c_{\mathsf n}\varrho^{\mathsf n},\qquad c_{\mathsf n}=\frac{1}{(\mathsf n-1)!}\Bigl(2^{\mathsf n-1}+\frac{2^{\mathsf n}}{\mathsf n}-2\mathsf n\Bigr),
\]
the coefficients of $\varrho^0$, $\varrho^1$ and $\varrho^2$ being zero. Here $c_3=\frac13$, and $c_{\mathsf n}>0$ for $\mathsf n\ge4$ because $2^{\mathsf n-1}\ge2\mathsf n$, so the series is positive for $\varrho>0$ and $d^\ast_{\mathrm M}>d^\ast_{\mathrm C}$. As $f\to0$, the expansions of (c) give $d^\ast_{\mathrm M}/d^\ast_{\mathrm C}\to(6f\xi/g)^{1/2}/(2(f\xi/g)^{1/2})=\sqrt6/2$; as $f\to g\xi$, $\bar f_{\mathrm M}\to2$ and $d^\ast_{\mathrm M}\to\infty$ while $d^\ast_{\mathrm C}$ tends to its finite value at $\bar f_{\mathrm C}=1$.
\end{proof}

\begin{proof}[Proof of \cref{prop:abs-one-round}]
A member with home $x_i$ can visit a point $x$ and return within a round exactly when $2\lvert x-x_i\rvert\le v_{\mathrm{tr}}T_{\mathrm{tr}}$, which identifies the daily reach. Two members meet at $x$ within a round exactly when $x$ lies in both reaches. If they meet, $\lvert x_i-x_j\rvert\le\lvert x_i-x\rvert+\lvert x-x_j\rvert\le v_{\mathrm{tr}}T_{\mathrm{tr}}$; conversely, if $\lvert x_i-x_j\rvert\le v_{\mathrm{tr}}T_{\mathrm{tr}}$, the midpoint of their homes lies in both reaches.
\end{proof}

\begin{proof}[Proof of \cref{cor:abs-one-round-extent}]
Positions differ by at most $\ell_\Pi$ in every coordinate along every coupling, so $\varsigma_d^2=\max_i\sum_ja^\Pi_{ij}(x_{j,d}-x_{i,d})^2\le\ell_\Pi^2a^\Pi_{\max}$ and $\xi_{2,d}^2=\varsigma_d^2/(2\varepsilon^\Pi)\le\ell_\Pi^2a^\Pi_{\max}/(2\varepsilon^\Pi)$; \cref{cor:col-lstar} gives $\Lstar_d(r_0)\le\xi_{2,d}/(r_0-1)^{1/2}$. By the coupling formulas of \cref{thm:col-diffusive-precision}, a coupling between agents $i$ and $j$ is nonzero in the reversible class only if some agent $k$ attends both, and in the normal class only if both attend some agent $k$, self weights included; by the triangle inequality through $k$ the physical distance between $i$ and $j$ is then at most $2\ell_{\mathrm{att}}$. The last bound uses \cref{prop:abs-one-round}.
\end{proof}

\begin{proof}[Proof of \cref{prop:abs-elasticity}]
(a) If $\Psi(x)=\Psi(1)x^{-\alpha_{\mathrm{el}}}$, then $\mathrm{el}_\Psi\equiv\alpha_{\mathrm{el}}$. Conversely, if for every $d$ the value $\mathrm{el}_\Psi(d/\xi)$ does not depend on $\xi>0$, then $\mathrm{el}_\Psi$ is constant on $(0,\infty)$, say equal to $\alpha_{\mathrm{el}}$, and integrating $\Psi'/\Psi=-\alpha_{\mathrm{el}}/x$ gives $\Psi(x)=\Psi(1)x^{-\alpha_{\mathrm{el}}}$.

(b) We use the representation $K_\nu(x)=\int_0^\infty e^{-x\cosh\mathsf s}\cosh(\nu\mathsf s)\mathrm d\mathsf s$ for $x>0$ \citep[eq.~10.32.9]{olver2010nist}. Differentiating under the integral gives $K_0'=-K_1$ and $K_0''=\int_0^\infty e^{-x\cosh\mathsf s}\cosh^2\mathsf s\mathrm d\mathsf s$. Since $\frac{\mathrm d}{\mathrm d\mathsf s}[-e^{-x\cosh\mathsf s}\sinh\mathsf s]=(x\sinh^2\mathsf s-\cosh\mathsf s)e^{-x\cosh\mathsf s}$ and the bracket vanishes at both ends, $xK_0''+K_0'-xK_0=\int_0^\infty(x\sinh^2\mathsf s-\cosh\mathsf s)e^{-x\cosh\mathsf s}\mathrm d\mathsf s=0$. With $\mathrm{el}_{K_0}=-xK_0'/K_0$,
\[
\mathrm{el}_{K_0}'=-\frac{K_0'+xK_0''}{K_0}+x\Bigl(\frac{K_0'}{K_0}\Bigr)^2=-x+\frac{\mathrm{el}_{K_0}^2}x=\frac{\mathrm{el}_{K_0}^2-x^2}x.
\]
Since $K_1(x)-K_0(x)=\int_0^\infty(\cosh\mathsf s-1)e^{-x\cosh\mathsf s}\mathrm d\mathsf s>0$, we have $\mathrm{el}_{K_0}(x)>x>0$, hence $\mathrm{el}_{K_0}'>0$. As $x\to0$, the substitution $\mathsf v=x\sinh\mathsf s$ gives $xK_1(x)=\int_0^\infty e^{-(x^2+\mathsf v^2)^{1/2}}\mathrm d\mathsf v\to1$ by monotone convergence, while $K_0(x)\ge e^{-1}\arccosh(1/x)\to\infty$, because the integrand is at least $e^{-1}$ where $\cosh\mathsf s\le1/x$; so $\mathrm{el}_{K_0}(x)\to0$. As $x\to\infty$, the expansion $K_\nu(x)=(\pi/2x)^{1/2}e^{-x}[1+(4\nu^2-1)/(8x)+O(x^{-2})]$ \citep[eqs.~10.40.2 and 10.17.1]{olver2010nist} gives $K_1/K_0=1+1/(2x)+O(x^{-2})$, so $\mathrm{el}_{K_0}(x)-x\to\frac12$. Finally $\mathrm{el}(d)=\mathrm{el}_{K_0}(d/\xi)$, and $d/\xi$ decreases as $\xi$ increases.

(c) The substitution $\varrho=d/\xi$ turns $\int_0^\infty e^{-d/\xi}\xi^{-1-\alpha_{\mathrm{el}}}\mathrm d\xi$ into $d^{-\alpha_{\mathrm{el}}}\int_0^\infty\varrho^{\alpha_{\mathrm{el}}-1}e^{-\varrho}\mathrm d\varrho$ for every $d>0$, which is finite for $\alpha_{\mathrm{el}}>0$; the weight is not integrable at $\xi=0$, and the profile diverges as $d\to0$.
\end{proof}

\begin{proof}[Proof of \cref{prop:abs-hierarchy}]
(a) Let $\mathsf B_{l}(i,i')=b^{-l}\one[i,i'\text{ lie in the same level-}l\text{ block}]$ for $0\le l\le n_{\mathrm H}$, the orthogonal projector onto functions constant on level-$l$ blocks, with $\mathsf B_0=I$. For $1\le l\le n_{\mathrm H}$, the range of $\mathsf B_{l-1}-\mathsf B_l$ consists of functions constant on level-$(l-1)$ blocks with zero sum on every level-$l$ block, and it is spanned by such functions supported on single level-$l$ blocks. Let $f$ be one of them, supported on the block $\mathfrak b$. For $i\in\mathfrak b$, the cousins of $i$ at levels below $l$ lie in $i$'s level-$(l-1)$ block and share the value $f_i$; the cousins at level $l$ are the rest of $\mathfrak b$ and carry in total $0-b^{l-1}f_i$; and the cousins at higher levels lie outside $\mathfrak b$. Hence
\[
(Lf)_i=\sum_{l'=1}^{n_{\mathrm H}}A_{\mathrm h}\mathfrak q^{l'-1}f_i-\sum_{l'<l}A_{\mathrm h}\mathfrak q^{l'-1}f_i+a^{\mathrm H}_lb^{l-1}f_i=\mathsf g_lf_i,
\]
using $a^{\mathrm H}_lb^{l-1}=A_{\mathrm h}\mathfrak q^{l-1}/(b-1)$. For $i\notin\mathfrak b$, all members of $\mathfrak b$ are cousins of $i$ at one level, so $(Lf)_i$ is a multiple of $\sum_{\mathfrak b}f=0$. Constants have eigenvalue 0, and $\mathsf B_{n_{\mathrm H}}+\sum_{l=1}^{n_{\mathrm H}}(\mathsf B_{l-1}-\mathsf B_l)=I$, so
\[
\mathsf G=\frac{\mathsf B_{n_{\mathrm H}}}{\varepsilon^\Pi}+\sum_{l=1}^{n_{\mathrm H}}\frac{\mathsf B_{l-1}-\mathsf B_l}{\varepsilon^\Pi+\mathsf g_l}.
\]
For a pair at level $l\ge1$, $\mathsf B_{l''}(i,i')=b^{-l''}$ for $l''\ge l$ and $0$ for $l''<l$, so $(\mathsf B_{l'-1}-\mathsf B_{l'})(i,i')$ is $(b-1)b^{-l'}$ for $l'>l$, $-b^{-l}$ for $l'=l$ and $0$ for $l'<l$; on the diagonal it is $(b-1)b^{-l'}$ for every $l'\ge1$. This gives (a).

(b) As $n_{\mathrm H}\to\infty$ at fixed $\varepsilon^\Pi>0$, the constant term $b^{-n_{\mathrm H}}/\varepsilon^\Pi$ tends to zero, $\mathsf g_l$ increases to $\mathsf c_\infty A_{\mathrm h}\mathfrak q^{l-1}$ with $\mathsf c_\infty=1/(1-\mathfrak q)+1/(b-1)$ (here $\mathfrak q<1$ is used), and the terms $(b-1)b^{-l'}/(\varepsilon^\Pi+\mathsf g_{l'})$ are dominated by the summable $(b-1)b^{-l'}/\varepsilon^\Pi$, so the sums converge to their infinite versions. As $\varepsilon^\Pi\to0$ these increase to sums with terms $(b-1)b^{-l'}/(\mathsf c_\infty A_{\mathrm h}\mathfrak q^{l'-1})=(b-1)\mathfrak q(b\mathfrak q)^{-l'}/(\mathsf c_\infty A_{\mathrm h})$, which are summable because $b\mathfrak q>1$; by monotone convergence the limits are
\[
\mathsf G(l)\to\frac{\mathfrak q}{\mathsf c_\infty A_{\mathrm h}}\Bigl[(b-1)\sum_{l'>l}(b\mathfrak q)^{-l'}-(b\mathfrak q)^{-l}\Bigr]=\frac{\mathfrak q}{\mathsf c_\infty A_{\mathrm h}}\frac{b(1-\mathfrak q)}{b\mathfrak q-1}(b\mathfrak q)^{-l},\qquad\mathsf G(0)\to\frac{\mathfrak q}{\mathsf c_\infty A_{\mathrm h}}\frac{b-1}{b\mathfrak q-1},
\]
and their ratio is the stated limit.

(c) With $l-1=D_e\ln(d_l/d_1)/\ln b$, $(b\mathfrak q)^{-l}=(b\mathfrak q)^{-1}(d_l/d_1)^{-D_e\ln(b\mathfrak q)/\ln b}$. The limit in (b) does not involve $A_{\mathrm h}$, and a change of the unit of distance multiplies $d_l$ and $d_1$ by the same factor.
\end{proof}

\begin{proof}[Proof of \cref{prop:abs-ultrametric-cap}]
For scalar errors the unconditional correlation is
\[
\Gamma_{ij}=\frac{\Sigma_{\mathrm{cm}}+\Sigma^c_{ij}}{\bigl((\Sigma_{\mathrm{cm}}+\Sigma^c_{ii})(\Sigma_{\mathrm{cm}}+\Sigma^c_{jj})\bigr)^{1/2}}\ge\frac{\Sigma_{\mathrm{cm}}}{\Sigma_{\mathrm{cm}}+\max_k\Sigma^c_{kk}},
\]
because $\Sigma^c_{ij}\ge0$. Every weight $-\ln\Gamma_{ij}$ is therefore at most $\ln(1+\max_k\Sigma^c_{kk}/\Sigma_{\mathrm{cm}})$. For a nonempty $C\ne V$ this bounds $\varpi_{\mathrm{out}}(C)$, a minimum over a nonempty set of pairs, while $u_{\mathrm{in}}(C)\ge0$. For the level-0 network, \cref{lem:fw-common-mode} splits the error into the common response to the world innovations and a part driven by local noise, independent of each other; with homogeneous row sums $W\one=\lambda\one$, the common part is $-\sum_{n\ge0}\lambda^nw_{t-1-n}\one$, with variance $Q/(1-\lambda^2)$, and $W\ge0$ makes $\Sigma^c=\sum_nW^n\diag(q_i)(W^\top)^n$ entrywise nonnegative for local noise variances $q_i$. The cap equals $\ln(1+(1-\lambda^2)\max_k\Sigma^c_{kk}/Q)$. For $W=\lambda P$ with $P$ irreducible, aperiodic and stochastic, $P^n=\one\pi^\top+R_n$ with $\lVert R_n\rVert\le c\rho^n$ for some $c$ and $\rho<1$ that depend on $P$ alone, so
\[
\Sigma^c=\sum_{n\ge0}\lambda^{2n}P^nQ^\zeta(P^\top)^n=\frac{\pi^\top Q^\zeta\pi}{1-\lambda^2}\one\one^\top+B_\lambda,\qquad B_\lambda=\sum_{n\ge0}\lambda^{2n}\bigl(\one\pi^\top Q^\zeta R_n^\top+R_nQ^\zeta\pi\one^\top+R_nQ^\zeta R_n^\top\bigr),
\]
and $\lVert B_\lambda\rVert$ is bounded uniformly in $\lambda\le1$ by a convergent geometric series. Hence $(1-\lambda^2)\Sigma^c_{kk}\to\pi^\top Q^\zeta\pi$ for every $k$, and the cap tends to $\ln(1+\pi^\top Q^\zeta\pi/Q)$. For symmetric $P$ the stationary vector is $\one/N$, and $Q^\zeta=qI$ gives $\pi^\top Q^\zeta\pi=q/N$. Aperiodicity is needed: for the swap of two agents with $Q^\zeta=\diag(q_1,q_2)$ the limit is $(q_1+q_2)/2$ against $\pi^\top Q^\zeta\pi=(q_1+q_2)/4$. For the bound at every $\lambda$, let $P$ be symmetric and doubly stochastic and $Q^\zeta=qI$, and write $P=\sum_k\mu_kv_kv_k^\top$ with orthonormal eigenvectors, $v_1=\one/\sqrt N$, $\mu_1=1$ and $\lvert\mu_k\rvert\le1$. Then $\Sigma^c=q\sum_{n\ge0}\lambda^{2n}P^{2n}=q\sum_k(1-\lambda^2\mu_k^2)^{-1}v_kv_k^\top$, every term has a nonnegative diagonal, and the term $k=1$ contributes $q/(N(1-\lambda^2))$ to every diagonal entry, so $(1-\lambda^2)\Sigma^c_{ii}\ge q/N$ for every $i$ and every $\lambda\in(0,1)$; neither irreducibility nor aperiodicity is used.
\end{proof}
